% !TEX root = ../main.tex
\corrigesChapitre{10}{Variables aléatoires}
\corrigesSouples

\begin{blocLivre}{cor-10-1}
\begin{corrige}{1}{Une variable indicatrice}
\demarche{Regrouper les issues qui donnent la même valeur de la variable.}
\begin{questions}
\item Trois issues sont paires sur six : $P(X=1)=P(X=0)=\res{1/2}$.
\item $E(X)=0\times1/2+1\times1/2=\res{1/2}$. $P(X\leq0)=P(X=0)=\res{1/2}$.
\end{questions}
\label{figlivre-r10-01-loi-cor}
\figureLivre[corrige]{r10-01-loi}

\end{corrige}
\end{blocLivre}
\begin{blocLivre}{cor-10-2}
\begin{corrige}{2}{Lire et exploiter une loi}
\demarche{Sélectionner les valeurs autorisées par chaque évènement avant d’additionner leurs probabilités.}
\begin{questions}
\item On additionne les masses concernées : $\res{0{,}7}$ et $\res{0{,}3}$.
\item $E(X)=-0{,}2+0{,}6=\res{0{,}4}$. $E(X^2)=0{,}2+4\times0{,}3=\res{1{,}4}$.
\end{questions}
\end{corrige}
\end{blocLivre}
\begin{blocLivre}{cor-10-3}
\begin{corrige}{3}{Compléter une loi}
\demarche{Compléter les probabilités à 1 avant de calculer les moments.}
\begin{questions}
\item La somme vaut 1, donc $p=1-0{,}1-0{,}2-0{,}4=\res{0{,}3}$.
\item $E(X)=0{,}2+0{,}6+1{,}2=\res{2}$. $E(X^2)=0{,}2+1{,}2+3{,}6=5$. Par König-Huygens, $V(X)=5-2^2=\res{1}$.
\end{questions}
\label{figlivre-r10-03-loi-cor}
\figureLivre[corrige]{r10-03-loi}

\end{corrige}
\end{blocLivre}
\begin{blocLivre}{cor-10-4}
\begin{corrige}{4}{Une loi de Bernoulli}
\demarche{Utiliser les deux valeurs possibles dans les définitions de l’espérance et de la variance.}
\begin{questions}
\item À partir des deux valeurs possibles,
$E(X)=0\times0{,}7+1\times0{,}3=\res{0{,}3}$.
La définition de la variance donne
\[V(X)=(0-0{,}3)^2\times0{,}7+(1-0{,}3)^2\times0{,}3
=0{,}063+0{,}147=\res{0{,}21}.\]
On peut contrôler avec $E(X^2)-E(X)^2=0{,}3-0{,}09$.
\item L’écart type est la racine positive de la variance :
$\sigma(X)=\sqrt{0{,}21}=\res{\sqrt{21}/10\approx0{,}458}$.
\end{questions}

\end{corrige}
\end{blocLivre}
\begin{blocLivre}{cor-10-5}
\begin{corrige}{5}{Un jeu équitable}
\demarche{Le gain net inclut la perte éventuelle ; un jeu équitable a une espérance nulle.}
\begin{questions}
\item $P(G=5)=1/6$ et $P(G=-1)=5/6$. $E(G)=5/6-5/6=\res{0}$.
\item $E(G^2)=25/6+5/6=5$, donc $\res{V(G)=5}$. Le gain moyen théorique est nul ; cela ne garantit ni gain ni perte nulle sur une partie.
\end{questions}
\label{figlivre-r10-05-loi-cor}
\figureLivre[corrige]{r10-05-loi}

\end{corrige}
\end{blocLivre}
\begin{blocLivre}{cor-10-6}
\begin{corrige}{6}{Calculer la dispersion de deux façons}
\demarche{Comparer les deux expressions de la variance avec des fractions exactes.}
\begin{questions}
\item $E(X)=(-1)/5+0/2+2(3/10)=2/5$ et $E(X^2)=1/5+4(3/10)=7/5$.
Par König-Huygens, $V(X)=7/5-(2/5)^2=35/25-4/25=\res{31/25}$.
\item Les trois écarts à $2/5$ sont $-7/5$, $-2/5$ et $8/5$.
Leurs carrés pondérés donnent
\[V(X)=\frac{49}{25}\frac15+\frac4{25}\frac12+\frac{64}{25}\frac3{10}
=\frac{49+10+96}{125}=\frac{31}{25}.\]
Les deux méthodes donnent bien la même variance.
\item $\sigma(X)=\sqrt{31/25}=\res{\sqrt{31}/5}$.
La racine porte sur toute la variance.
\end{questions}

\end{corrige}
\end{blocLivre}
\begin{blocLivre}{cor-10-7}
\begin{corrige}{7}{Un changement d'unité}
\demarche{Transformer les valeurs et appliquer les règles de l’espérance et de la variance.}
\begin{questions}
\item La linéarité donne $E(Y)=3E(X)-2=3(2/5)-2=\res{-4/5}$.
La variance est multipliée par le carré de 3 :
$V(Y)=3^2V(X)=9(31/25)=\res{279/25}$.
Donc $\sigma(Y)=\sqrt{279/25}=\res{3\sqrt{31}/5}$.
La translation de $-2$ ne modifie pas la dispersion.
\item Les images de $-1$, 0 et 2 par $x\mapsto3x-2$ sont $-5$, $-2$ et 4.
Ces trois valeurs restent distinctes : chacune garde la probabilité de son antécédent, comme dans le tableau.
\end{questions}

\label{figlivre-r10-07-loi-cor}
\figureLivre[corrige]{r10-07-loi}

\end{corrige}
\end{blocLivre}
\begin{blocLivre}{cor-10-8}
\begin{corrige}{8}{Même moyenne, risques différents}
\demarche{Comparer les montants autour de leur moyenne, puis les probabilités de gain.}
\begin{questions}
\item Les deux espérances valent 0. Pour A, $V(A)=4$ et $\sigma(A)=2$ ; pour B, $V(B)=1$ et $\sigma(B)=1$.
\item Les deux probabilités sont $\res{1/2}$.
\item A présente des écarts plus grands à la moyenne, sans avoir une plus grande probabilité de gain. La variance mesure la dispersion des montants.
\end{questions}
\end{corrige}
\end{blocLivre}
\begin{blocLivre}{cor-10-9}
\begin{corrige}{9}{Linéarité sans indépendance}
\demarche{La linéarité de l’espérance reste valable lorsque les variables dépendent l’une de l’autre.}
\begin{questions}
\item $E(X)=(1+\cdots+6)/6=\res{7/2}$. Par linéarité $E(Y)=7-E(X)=\res{7/2}$.
\item $X+Y=7$ pour toute issue : $E(X+Y)=7$ et $V(X+Y)=0$. On retrouve $E(X)+E(Y)=7$.
\item Les deux évènements coïncident : leur intersection a probabilité $1/6$, différente de $(1/6)^2=1/36$. La linéarité a donc été utilisée sans indépendance.
\end{questions}
\end{corrige}
\end{blocLivre}
\begin{blocLivre}{cor-10-10}
\begin{corrige}{10}{Deux variables sur le même univers}
\demarche{Suivre les deux variables sur les mêmes quatre issues avant de les combiner.}
\begin{questions}
\item Les deux pièces indépendantes donnent quatre issues équiprobables, notées dans l’ordre des lancers.
Pour chaque ligne du tableau, on compte les faces pour $X$, on lit la première pièce pour $Y$, puis on calcule $Z=X-2Y$.
Par exemple, sur $FP$, $X=1$ et $Y=1$, donc $Z=-1$.
\item $E(X)=1$ et $E(Y)=1/2$. La linéarité ne demande pas d’indépendance entre $X$ et $Y$ :
$E(Z)=E(X)-2E(Y)=1-2(1/2)=\res{0}$.
\item La valeur $-1$ apparaît sur $FP$, la valeur 1 sur $PF$, et la valeur 0 sur $PP$ ou $FF$.
Donc $P(Z=-1)=1/4$, $P(Z=0)=1/2$, $P(Z=1)=1/4$.
On vérifie $E(Z)=(-1)/4+1/4=0$ et
$V(Z)=E(Z^2)-0^2=1/4+1/4=\res{1/2}$.
\end{questions}

\label{figlivre-r10-10-issues-cor}
\figureLivre[corrige]{r10-10-issues}

\end{corrige}
\end{blocLivre}
\begin{blocLivre}{cor-10-11}
\begin{corrige}{11}{Choisir une valeur de prévision}
\demarche{Développer le carré de l’erreur et utiliser la linéarité.}
\begin{questions}
\item $E(X)=-1/4+1/2+3/4=\res{1}$. $E(X^2)=1/4+1/2+9/4=\res{3}$.
\item Par développement et linéarité, $E((X-a)^2)=E(X^2)-2aE(X)+a^2=3-2a+a^2=(a-1)^2+2$.
\item Le carré est minimal pour $\res{a=1}$ ; l’erreur minimale vaut $\res{2}$, qui est la variance de $X$.
\end{questions}
\end{corrige}
\end{blocLivre}
\begin{blocLivre}{cor-10-12}
\begin{corrige}{12}{Une loterie et son prix d'entrée}
\demarche{Distinguer le lot brut, le prix d’entrée et le gain net.}
\begin{questions}
\item Il reste $100-3-10=87$ tickets sans lot.
La loi de $L$ est donnée dans le tableau et son espérance vaut
$E(L)=0(87/100)+5(10/100)+20(3/100)=\res{1{,}10\text{ euro}}$.
\item Le prix est payé à chaque partie : $G=L-m$, donc $E(G)=E(L)-m=1{,}1-m$.
Un jeu équitable vérifie $E(G)=0$, ce qui donne $\res{m=1{,}10\text{ euro}}$.
\item Pour $m=2$, $E(G)=\res{-0{,}90\text{ euro}}$.
Soustraire une constante conserve la variance : $V(G)=V(L)$.
Comme $E(L^2)=0+25(0{,}10)+400(0{,}03)=14{,}5$,
$V(G)=14{,}5-1{,}1^2=\res{13{,}29}$ euros carrés.
\item Un gain net positif impose $L>2$ : seuls les lots de 5 et 20 euros conviennent.
La probabilité est donc $10/100+3/100=\res{13/100}$.
\end{questions}

\label{figlivre-r10-12-loi-cor}
\figureLivre[corrige]{r10-12-loi}

\end{corrige}
\end{blocLivre}
\begin{blocLivre}{cor-10-13}
\begin{corrige}{13}{Un jeu sur trois pièces}
\demarche{Compter les huit issues équiprobables et les regrouper par nombre de faces.}
\begin{questions}
\item Les huit issues ont chacune la probabilité $1/8$.
Le tableau les regroupe par nombre de faces : on trouve respectivement 1, 3, 3 et 1 issues,
donc les probabilités $1/8$, $3/8$, $3/8$, $1/8$ pour $X=0,1,2,3$.
\item $E(X)=[0\times1+1\times3+2\times3+3\times1]/8=\res{3/2}$.
De même $E(X^2)=(3+12+9)/8=3$, donc $V(X)=3-(3/2)^2=\res{3/4}$.
\item Le gain net est $G=4X-5$ : la mise est retranchée quel que soit le résultat.
Ainsi $E(G)=4(3/2)-5=\res{1\text{ euro}}$ et $V(G)=4^2(3/4)=\res{12}$ euros carrés.
\item Pour une mise $m$, l’équité demande $4E(X)-m=0$, donc $\res{m=6\text{ euros}}$.
Avec la mise 5, $4X-5>0\iff X>5/4$. Comme $X$ est entier, cela équivaut à $X\geq2$,
de probabilité $(3+1)/8=\res{1/2}$.
\end{questions}

\label{figlivre-r10-13-issues-cor}
\figureLivre[corrige]{r10-13-issues}

\end{corrige}
\end{blocLivre}
\begin{blocLivre}{cor-10-14}
\begin{corrige}{14}{Un modèle de garantie}
\demarche{Retirer le coût de réparation du prix reçu, puis calculer la dispersion.}
\begin{questions}
\item Sans réparation, le résultat net vaut 15 euros.
Avec réparation, il vaut $15-500=-485$ euros.
La loi du tableau donne $E(G)=15(0{,}98)-485(0{,}02)=\res{5\text{ euros}}$.
\item Posons $I=1$ si une réparation est nécessaire, et $I=0$ sinon.
Alors $G=15-500I$, avec $V(I)=0{,}02\times0{,}98$.
Donc $V(G)=500^2\times0{,}02\times0{,}98=\res{4900}$ euros carrés et
$\sigma(G)=\sqrt{4900}=\res{70\text{ euros}}$.
La dépense rare de 500 euros explique cette forte dispersion autour d’une espérance de 5 euros.
\item Si $G_1,\ldots,G_{100}$ sont les résultats des garanties,
$E(G_1+\cdots+G_{100})=100\times5=\res{500\text{ euros}}$.
Cette addition des espérances ne demande pas d’indépendance.
\end{questions}

\label{figlivre-r10-14-loi-cor}
\figureLivre[corrige]{r10-14-loi}

\end{corrige}
\end{blocLivre}
\begin{blocLivre}{cor-10-15}
\begin{corrige}{15}{Une aire aléatoire}
\demarche{Exprimer l’aire en fonction de la hauteur avant d’établir sa loi.}
\begin{questions}
\item Le triangle est rectangle en $A$, avec base 1 et hauteur $H$.
Ainsi $S=1\times H/2=H/2$.
Les trois valeurs possibles sont $1/2$, 1 et 2, avec les probabilités du tableau.
\item En utilisant directement la loi de $S$,
$E(S)=(1/2)(1/2)+1(1/4)+2(1/4)=\res{1}$.
Puis $E(S^2)=(1/4)(1/2)+1(1/4)+4(1/4)=11/8$.
Donc $V(S)=11/8-1^2=\res{3/8}$ et $P(S\geq1)=1/4+1/4=\res{1/2}$.
\item Développer le carré donne
$E((S-a)^2)=E(S^2)-2aE(S)+a^2=11/8-2a+a^2=\res{(a-1)^2+3/8}$.
\item Le carré $(a-1)^2$ est minimal en $a=1$.
La meilleure aire fixe est donc $\res{1}$ unité d’aire, pour une erreur quadratique moyenne de $3/8$.
\end{questions}

\label{figlivre-r10-15-loi-cor}
\figureLivre[corrige]{r10-15-loi}

\label{figlivre-r10-15-triangle-cor}
\figureLivre[corrige]{r10-15-triangle}

\end{corrige}
\end{blocLivre}
